\section*{Conclusiones}


La calibración de la termocupla permitió establecer una relación lineal clara entre la temperatura y el voltaje de salida del sistema de acondicionamiento, evidenciada por la regresión obtenida a partir de los datos experimentales. La ecuación de calibración demuestra que, dentro del rango de operación considerado (20 °C a 40 °C), el comportamiento del sistema es aproximadamente lineal, lo cual valida tanto el modelo teórico del termopar como el diseño del amplificador. Esta linealidad facilita la conversión directa de voltaje a temperatura y confirma que el circuito implementado es adecuado para aplicaciones de medición de temperatura en rangos bajos con buena repetibilidad.

La correcta supresión del offset en los amplificadores operacionales fue un factor determinante para lograr una calibración confiable, ya que la presencia de offset se ve amplificada proporcionalmente con la ganancia total del circuito. El uso de potenciómetros de ajuste en el LM741 permitió reducir el voltaje de offset a valores cercanos a cero en cada etapa, minimizando errores sistemáticos en la salida. Esto es especialmente relevante en la primera etapa, donde la ganancia es mayor y cualquier desbalance interno del operacional tiene un impacto más significativo en la medición final de la termocupla.

El filtrado mediante los capacitores $C_1$ y $C_2$, junto con la reducción del rizado en las fuentes de alimentación, contribuyó a estabilizar la señal de salida y a mejorar la calidad de la calibración experimental. Al limitar el ancho de banda a una frecuencia de corte de 6 Hz y suprimir el ruido de la fuente DC, se logró una señal más limpia y estable, lo que permitió obtener mediciones de voltaje más consistentes para cada punto de temperatura. En conjunto, el diseño del filtrado, la ganancia total del sistema y la calibración experimental aseguran un desempeño adecuado del circuito para la medición precisa de temperatura mediante termocupla.
